\documentclass[a4paper,11pt]{article}

% ── Packages ──────────────────────────────────────────────────
\usepackage[a4paper, top=2.5cm, bottom=2.5cm, left=2cm, right=2cm]{geometry}
\usepackage[T1]{fontenc}
\usepackage[utf8]{inputenc}
\usepackage{lmodern}
\usepackage{microtype}
\usepackage{xcolor}
\usepackage{booktabs}
\usepackage{tabularx}
\usepackage{array}
\usepackage{listings}
\usepackage{parskip}
\usepackage{titlesec}
\usepackage{hyperref}
\usepackage{needspace}

% ── Colors ────────────────────────────────────────────────────
\definecolor{accent}{HTML}{1a3c6e}
\definecolor{lightbg}{HTML}{e8eef7}
\definecolor{codebg}{HTML}{f4f4f4}
\definecolor{tablegray}{HTML}{aaaaaa}
\definecolor{gray}{HTML}{555555}

% ── Hyperref setup ────────────────────────────────────────────
\hypersetup{
    colorlinks=true,
    linkcolor=accent,
    urlcolor=accent,
    pdftitle={MOD-08 Integration \& Test Module},
    pdfauthor={Ayse Feyza SERBEST, Gorkem UYSAL}
}

% ── Section title formatting ──────────────────────────────────
\titleformat{\section}
  {\large\bfseries\color{accent}}
  {\thesection.}{0.5em}{}
  [\vspace{-0.4em}\textcolor{lightbg}{\rule{\linewidth}{0.5pt}}]

\titleformat{\subsection}
  {\normalsize\bfseries\color{accent}}
  {\thesubsection}{0.5em}{}

\titlespacing{\section}{0pt}{18pt}{6pt}
\titlespacing{\subsection}{0pt}{12pt}{4pt}

% ── Code listing style ────────────────────────────────────────
\lstset{
    backgroundcolor=\color{codebg},
    basicstyle=\ttfamily\small,
    breaklines=true,
    frame=single,
    framerule=0pt,
    rulecolor=\color{codebg},
    xleftmargin=8pt,
    xrightmargin=8pt,
    aboveskip=6pt,
    belowskip=8pt,
    language=C
}

% ── Table helpers ─────────────────────────────────────────────
\newcolumntype{L}[1]{>{\raggedright\arraybackslash}p{#1}}
\newcolumntype{C}[1]{>{\centering\arraybackslash}p{#1}}

\newcommand{\theadbg}{\rowcolor{accent}}
\newcommand{\theadtext}[1]{\textcolor{white}{\textbf{#1}}}

\arrayrulecolor{tablegray}

% ── Document ──────────────────────────────────────────────────
\begin{document}

% ── Cover ─────────────────────────────────────────────────────
\textcolor{gray}{\small CSE 396 --- Computer Engineering Project}\\
\textcolor{gray}{\small Autonomous Fire Detection and Suppression Rover}

\vspace{8pt}
{\color{accent}\rule{\linewidth}{2pt}}
\vspace{6pt}

{\LARGE\bfseries\color{accent} MOD-08: Integration \& Test Module}\\[4pt]
{\large Technical Module Documentation}

\vspace{4pt}
{\color{gray}\rule{\linewidth}{0.5pt}}
\vspace{10pt}

\begin{tabular}{
  >{\columncolor{accent}}L{6.46cm}
  >{\columncolor{accent}}L{4.25cm}
  >{\columncolor{accent}}L{5.3cm}
}
\theadtext{Name} & \theadtext{Student ID} & \theadtext{Assigned Modules} \\
\end{tabular}

\vspace{0pt}
\begin{tabular}{L{6.46cm} L{4.25cm} L{5.3cm}}
\rowcolor{white}   Ayşe Feyza SERBEST & 220104004052 & MOD-07, MOD-08 \\
\rowcolor{lightbg} Görkem UYSAL        & 230104004174 & MOD-07, MOD-08 \\
\end{tabular}

\vspace{16pt}

% ── 1. Purpose ────────────────────────────────────────────────
\needspace{6\baselineskip}
\section{Purpose}

MOD-08 serves as the system-wide validation and quality assurance layer for the Autonomous
Fire Detection and Suppression Rover. Its primary responsibility is to verify that all
eight modules operate correctly both individually and in concert before and during field
deployment.

The module provides two complementary testing modes: a \textbf{real hardware self-check}
that exercises every physical peripheral on the rover, and a \textbf{simulation mode} that
injects synthetic sensor data into the system, allowing full Finite State Machine (FSM)
validation without requiring a physical fire source or any hardware.

% ── 2. Dependencies ───────────────────────────────────────────
\needspace{14\baselineskip}
\section{Dependencies}

\begin{tabular}{L{4.76cm} L{2.38cm} L{9.86cm}}
\toprule
\rowcolor{accent}
\theadtext{Module Name} & \theadtext{ID} & \theadtext{Usage in MOD-08} \\
\midrule
\rowcolor{white}   Flame Sensing HW    & \texttt{MOD-01} & Sensor presence validation during self-check \\
\rowcolor{lightbg} Flame Detection SW  & \texttt{MOD-02} & Full detection pipeline verification via \texttt{flame\_data\_t} \\
\rowcolor{white}   Locomotion HW       & \texttt{MOD-03} & Motor driver response test \\
\rowcolor{lightbg} Motion Control SW   & \texttt{MOD-04} & Movement command verification via \texttt{motion\_direction\_t} \\
\rowcolor{white}   Distance Sensing    & \texttt{MOD-05} & Ultrasonic obstacle detection validation \\
\rowcolor{lightbg} Fire Suppression HW & \texttt{MOD-06} & Relay and pump activation test \\
\rowcolor{white}   Autonomy SW         & \texttt{MOD-07} & End-to-end FSM behavior testing with injected triggers \\
\bottomrule
\end{tabular}

% ── 3. API Summary ────────────────────────────────────────────
\needspace{20\baselineskip}
\section{API Summary}

\subsection{Public Functions}

\begin{tabular}{L{3.74cm} L{3.74cm} L{2.38cm} L{7.14cm}}
\toprule
\rowcolor{accent}
\theadtext{Function} & \theadtext{Parameters} & \theadtext{Return} & \theadtext{Description} \\
\midrule
\rowcolor{white}   \texttt{test\_init()}                  & \texttt{void}                                      & \texttt{void}            & Initialises internal state and links to all dependent modules. \\
\rowcolor{lightbg} \texttt{test\_set\_mode()}             & \texttt{test\_mode\_t mode}                        & \texttt{void}            & Switches between \texttt{TEST\_MODE\_REAL} and \texttt{TEST\_MODE\_SIMULATION}. \\
\rowcolor{white}   \texttt{test\_run\_self\_check()}      & \texttt{void}                                      & \texttt{test\_status\_t} & Sequentially validates all hardware peripherals; returns first failure code or \texttt{TEST\_OK}. \\
\rowcolor{lightbg} \texttt{test\_get\_report()}           & \texttt{test\_report\_t *report}                   & \texttt{void}            & Fills caller-owned struct with per-subsystem pass/fail flags. \\
\rowcolor{white}   \texttt{test\_set\_simulation\_data()} & \texttt{const test\_simulation\_input\_t *data}    & \texttt{void}            & Injects synthetic flame and distance data for simulation runs. \\
\rowcolor{lightbg} \texttt{test\_update()}                & \texttt{void}                                      & \texttt{void}            & Periodic monitoring tick; logs anomalies during continuous operation. \\
\bottomrule
\end{tabular}

\subsection{Data Types}

\begin{tabular}{L{4.25cm} L{2.04cm} L{10.71cm}}
\toprule
\rowcolor{accent}
\theadtext{Type} & \theadtext{Kind} & \theadtext{Fields / Values} \\
\midrule
\rowcolor{white}   \texttt{test\_status\_t}            & enum   & \texttt{TEST\_OK}, \texttt{TEST\_FAIL\_SENSOR}, \texttt{TEST\_FAIL\_DISTANCE}, \texttt{TEST\_FAIL\_MOTOR}, \texttt{TEST\_FAIL\_PUMP} \\
\rowcolor{lightbg} \texttt{test\_mode\_t}              & enum   & \texttt{TEST\_MODE\_REAL}, \texttt{TEST\_MODE\_SIMULATION} \\
\rowcolor{white}   \texttt{test\_report\_t}            & struct & \texttt{sensor\_ok}, \texttt{distance\_ok}, \texttt{motor\_ok}, \texttt{pump\_ok} (all \texttt{int}) \\
\rowcolor{lightbg} \texttt{test\_simulation\_input\_t} & struct & \texttt{flame\_data} (\texttt{flame\_data\_t} from MOD-02), \texttt{distance\_cm} (\texttt{uint16\_t}) \\
\bottomrule
\end{tabular}

% ── 4. Module Connections ─────────────────────────────────────
\needspace{16\baselineskip}
\section{Module Connections}

MOD-08 occupies a unique position in the architecture: it has \textbf{read access to every
module's output} for validation purposes, and \textbf{write/inject access} to MOD-02,
MOD-05, and MOD-07 for simulation. In production mode it is passive; in test mode it
becomes the data source.

\vspace{6pt}
\begin{tabular}{L{1.87cm} L{1.87cm} L{3.74cm} L{2.72cm} L{6.8cm}}
\toprule
\rowcolor{accent}
\theadtext{From} & \theadtext{To} & \theadtext{Signal / Data} & \theadtext{Direction} & \theadtext{Notes} \\
\midrule
\rowcolor{white}   \texttt{MOD-08} & \texttt{MOD-02} & \texttt{simulate\_flame}    & 08 $\rightarrow$ 02 & Injected \texttt{flame\_data\_t} for simulation runs \\
\rowcolor{lightbg} \texttt{MOD-08} & \texttt{MOD-05} & \texttt{simulate\_distance} & 08 $\rightarrow$ 05 & Injected \texttt{distance\_cm} for simulation runs \\
\rowcolor{white}   \texttt{MOD-08} & \texttt{MOD-07} & \texttt{test\_triggers}     & 08 $\rightarrow$ 07 & FSM state injection for behavioral testing \\
\rowcolor{lightbg} \texttt{MOD-01} & \texttt{MOD-08} & \texttt{sensor\_status}     & 01 $\rightarrow$ 08 & Sensor presence confirmed during self-check \\
\rowcolor{white}   \texttt{MOD-03} & \texttt{MOD-08} & \texttt{motor\_status}      & 03 $\rightarrow$ 08 & Motor driver response during self-check \\
\rowcolor{lightbg} \texttt{MOD-05} & \texttt{MOD-08} & \texttt{dist\_reading}      & 05 $\rightarrow$ 08 & Ultrasonic ping result for validation \\
\rowcolor{white}   \texttt{MOD-06} & \texttt{MOD-08} & \texttt{pump\_status}       & 06 $\rightarrow$ 08 & Relay activation confirmation \\
\rowcolor{lightbg} \texttt{MOD-07} & \texttt{MOD-08} & \texttt{fsm\_state}         & 07 $\rightarrow$ 08 & Current FSM state read for behavioral validation \\
\bottomrule
\end{tabular}

% ── 5. Quick-Start Integration ────────────────────────────────
\needspace{24\baselineskip}
\section{Quick-Start Integration}

\subsection{Real Hardware Self-Check}

Call \texttt{test\_run\_self\_check()} inside \texttt{setup()} before the main loop begins.
If any peripheral fails, the rover halts and the failure code can be read via
\texttt{test\_get\_report()}.

\begin{lstlisting}
#include "system_test.h"

void setup() {
    test_init();
    test_set_mode(TEST_MODE_REAL);

    if (test_run_self_check() != TEST_OK) {
        test_report_t report;
        test_get_report(&report);
        // Log report.sensor_ok, report.motor_ok, etc.
        while (1); // Halt on failure
    }
}

void loop() {
    test_update(); // Optional: continuous monitoring
}
\end{lstlisting}

\subsection{Simulation Mode --- FSM Testing Without Hardware}

Simulation mode allows developers to validate the full decision pipeline on a laptop or
in a CI environment. Synthetic \texttt{flame\_data\_t} and distance values are injected
directly into MOD-02 and MOD-05, exercising MOD-07's state transitions without any
physical sensors.

\begin{lstlisting}
#include "system_test.h"

void test_center_flame_approach() {
    test_simulation_input_t sim = {
        .flame_data = {
            .detected      = 1,
            .direction     = FLAME_DIR_CENTER,
            .max_intensity = 800
        },
        .distance_cm = 40
    };

    test_set_mode(TEST_MODE_SIMULATION);
    test_set_simulation_data(&sim);
    test_update();
    // Assert: autonomy_get_current_state() == STATE_APPROACH
}
\end{lstlisting}

% ── 6. Operational Flow ───────────────────────────────────────
\needspace{12\baselineskip}
\section{Operational Flow}

The table below summarises the two operational paths through MOD-08.

\vspace{6pt}
\begin{tabular}{C{0.5cm} L{7.99cm} L{7.99cm}}
\toprule
\rowcolor{accent}
\theadtext{\#} & \theadtext{Real Mode} & \theadtext{Simulation Mode} \\
\midrule
\rowcolor{white}   1 & \texttt{test\_init()} --- link all modules                  & \texttt{test\_init()} --- link all modules \\
\rowcolor{lightbg} 2 & \texttt{test\_set\_mode(TEST\_MODE\_REAL)}                   & \texttt{test\_set\_mode(TEST\_MODE\_SIMULATION)} \\
\rowcolor{white}   3 & \texttt{test\_run\_self\_check()} --- ping each peripheral   & \texttt{test\_set\_simulation\_data()} --- load scenario \\
\rowcolor{lightbg} 4 & \texttt{test\_get\_report()} --- read pass/fail flags        & \texttt{test\_update()} --- inject data into MOD-02/05/07 \\
\rowcolor{white}   5 & \texttt{test\_update()} --- continuous monitoring loop       & Read \texttt{autonomy\_get\_current\_state()} to assert \\
\bottomrule
\end{tabular}

% ── 7. Known Limitations & TODOs ─────────────────────────────
\needspace{18\baselineskip}
\section{Known Limitations \& TODOs}

\subsection{Current Limitations}

\begin{itemize}
  \item Test routines depend on controlled environment conditions; field noise may produce false failures.
  \item Simulation scenarios are limited to a single static input snapshot --- dynamic sequences are not yet supported.
  \item No detailed per-component error logging; \texttt{test\_report\_t} provides only binary pass/fail per subsystem.
  \item \texttt{test\_update()} does not yet implement continuous flame-loss or motor-stall detection.
\end{itemize}

\subsection{Planned Improvements (TODOs)}

\begin{itemize}
  \item[$\square$] Add granular error codes to \texttt{test\_report\_t} (e.g., which sensor index failed).
  \item[$\square$] Extend simulation to support multi-step scenarios with changing distance and flame direction over time.
  \item[$\square$] Implement automated test logging to serial output or SD card for field debugging.
  \item[$\square$] Integrate continuous system monitoring with alert thresholds in \texttt{test\_update()}.
  \item[$\square$] Add a dedicated stress-test mode that cycles \texttt{STATE\_EXTINGUISH} repeatedly to validate relay cooldown (MOD-06).
\end{itemize}

% ── 8. Version History ────────────────────────────────────────
\needspace{8\baselineskip}
\section{Version History}

\begin{tabular}{L{1.7cm} L{2.55cm} L{12.75cm}}
\toprule
\rowcolor{accent}
\theadtext{Version} & \theadtext{Date} & \theadtext{Changes} \\
\midrule
\rowcolor{white}   \texttt{v0.1} & 2026-03-28 & Initial draft --- \texttt{test\_init}, \texttt{test\_run\_self\_check}, \texttt{test\_simulate\_fire\_event}, basic self-check stubs. \\
\rowcolor{lightbg} \texttt{v0.2} & 2026-03-29 & Added \texttt{test\_mode\_t}, \texttt{test\_report\_t}, \texttt{test\_set\_simulation\_data}, \texttt{test\_update}; improved reporting and simulation control. \\
\bottomrule
\end{tabular}

\end{document}
